\documentclass[a4paper, 12pt]{article}
\usepackage[margin=2cm]{geometry}
\usepackage{times, amssymb, amsmath, graphicx, tabularx, booktabs, subfiles, amsthm,enumitem,xcolor,braket, subcaption} 

% Add this in the preamble of your LaTeX document

\usepackage{tikz}
\usetikzlibrary{arrows.meta}


\newcommand{\ans}[1]{\textcolor{red}{#1}}

\newcommand{\boxx}[1]{%
    \begin{tcolorbox}[colback=gray!10, colframe=gray!50, title=Where it is getting wrong?]
        #1
    \end{tcolorbox}%
}


\title{}
\author{}

\begin{document}
\maketitle
%------------Body-Starts--------------------


\textbf{1. Answer all questions of the following:}

\textbf{(a) What is meant by stopping potential in the photoelectric effect?}

Stopping potential is the minimum negative (retarding) potential applied to the collector plate which just stops the most energetic photoelectrons, making the photocurrent zero.  
If $K_{\text{max}}$ is the maximum kinetic energy of emitted electrons, then

\[
eV_0 = K_{\text{max}}
\]

\textbf{\textcolor{brown}{\textbf{Stopping potential is the minimum retarding potential required to reduce the photocurrent to zero by stopping the most energetic photoelectrons.}}}

\vspace{0.5cm}

\textbf{(b) What does the Davisson and Germer experiment demonstrate?}

The Davisson and Germer experiment demonstrated the diffraction of electrons by a crystal lattice, confirming the wave nature of electrons as proposed by de Broglie.

\textbf{\textcolor{brown}{\textbf{The experiment demonstrates the wave nature of electrons and verifies de Broglie’s hypothesis.}}}

\vspace{0.5cm}

\textbf{(c) Write down the expression for radius of a nucleus.}

The radius of a nucleus is given by

\[
R = R_0 A^{1/3}
\]

where $R_0 \approx 1.2 \times 10^{-15}\,\text{m}$ and $A$ is the mass number.

\textbf{\textcolor{brown}{\textbf{The radius of a nucleus is given by $R = R_0 A^{1/3}$.}}}

\vspace{0.5cm}

\textbf{(d) How does the volume term in the semi-empirical mass formula contribute to binding energy?}

In the semi-empirical mass formula, the volume term is proportional to the mass number $A$:

\[
E_v = a_v A
\]

Since each nucleon interacts with a nearly constant number of neighboring nucleons, this term contributes positively and increases linearly with $A$, thereby increasing the binding energy.

\textbf{\textcolor{brown}{\textbf{The volume term contributes positively to binding energy and is proportional to $A$, i.e., $E_v = a_v A$.}}}

\vspace{0.5cm}

\textbf{(e) Under what condition does population inversion take place in a laser medium?}

Population inversion occurs when the number of atoms in the excited state exceeds the number in the lower energy state:

\[
N_2 > N_1
\]

\textbf{\textcolor{brown}{\textbf{Population inversion occurs when the population of the upper energy level exceeds that of the lower level ($N_2 > N_1$).}}}




\textbf{2. A particle is described by the one-dimensional wave function}
\[
\psi(x) = A e^{-\alpha x^{2}}, \qquad -\infty < x < \infty
\]
\textbf{where $A$ is a positive constant. Write down the normalization condition for the wave function. Hence determine the value of the normalization constant $A$ in terms of $\alpha$.}

\vspace{0.3cm}

The normalization condition for a wave function is

\[
\int_{-\infty}^{\infty} |\psi(x)|^{2} \, dx = 1
\]

Since $\psi(x) = A e^{-\alpha x^{2}}$ is real,

\[
|\psi(x)|^{2} = A^{2} e^{-2\alpha x^{2}}
\]

Substituting into the normalization condition,

\[
\int_{-\infty}^{\infty} A^{2} e^{-2\alpha x^{2}} \, dx = 1
\]

Taking $A^{2}$ outside the integral,

\[
A^{2} \int_{-\infty}^{\infty} e^{-2\alpha x^{2}} \, dx = 1
\]

Using the Gaussian integral,

\[
\int_{-\infty}^{\infty} e^{-a x^{2}} \, dx = \sqrt{\frac{\pi}{a}}, \qquad a > 0
\]

Here $a = 2\alpha$, therefore

\[
\int_{-\infty}^{\infty} e^{-2\alpha x^{2}} \, dx
=
\sqrt{\frac{\pi}{2\alpha}}
\]

Substituting,

\[
A^{2} \sqrt{\frac{\pi}{2\alpha}} = 1
\]

Solving for $A$,

\[
A^{2} = \sqrt{\frac{2\alpha}{\pi}}
\]

\[
A = \left( \frac{2\alpha}{\pi} \right)^{1/4}
\]

\[
\textcolor{brown}{\mathbf{A = \left( \dfrac{2\alpha}{\pi} \right)^{1/4}}}
\]




\textbf{3. Differentiate between spontaneous emission and stimulated emission highlighting the role of an incident photon in each process.}

\vspace{0.3cm}

Consider an atom having two energy levels $E_2$ (excited state) and $E_1$ (lower state) such that

\[
E_2 > E_1
\]

If the atom is initially in the excited state, it can undergo transition to the lower state by emitting radiation of frequency

\[
h\nu = E_2 - E_1
\]

\textbf{Spontaneous Emission}

In spontaneous emission, the atom in the excited state $E_2$ returns to the lower state $E_1$ without the influence of any external radiation field. The emission occurs randomly due to internal atomic instability.

The emitted photon:

\[
h\nu = E_2 - E_1
\]

is emitted in a random direction with random phase and polarization.

The probability per unit time of spontaneous transition is given by Einstein’s coefficient $A_{21}$.

There is \emph{no incident photon} involved in initiating the process.

\vspace{0.4cm}

\textbf{Stimulated Emission}

In stimulated emission, an external photon of energy $h\nu$ such that

\[
h\nu = E_2 - E_1
\]

interacts with an atom already in the excited state $E_2$.

The incident photon induces the atom to return to the lower state $E_1$, resulting in the emission of a second photon.

The emitted photon has:

\[
\text{same frequency}, \quad \text{same phase}, \quad \text{same direction}, \quad \text{same polarization}
\]

as the incident photon.

Thus, two coherent photons emerge from the interaction.

The probability per unit time of stimulated transition is proportional to the radiation energy density $\rho(\nu)$ and is given by Einstein’s coefficient $B_{21}$.

\vspace{0.4cm}

\textbf{Key Differences}

\begin{center}
\begin{tabular}{|c|c|c|}
\hline
\textbf{Feature} & \textbf{Spontaneous Emission} & \textbf{Stimulated Emission} \\
\hline
Initiation & No incident photon & Requires incident photon \\
\hline
Direction of emitted photon & Random & Same as incident photon \\
\hline
Phase relation & Random phase & Same phase (coherent) \\
\hline
Rate dependence & Independent of radiation field & Proportional to radiation energy density $\rho(\nu)$ \\
\hline
\end{tabular}
\end{center}

\[
\textcolor{brown}{\mathbf{\text{Spontaneous emission occurs without any incident photon and produces incoherent radiation, whereas stimulated emission is induced by an incident photon and produces coherent radiation identical to the incident photon.}}}
\]


\textbf{4. A particle of mass $m$ is confined in a one-dimensional infinite square potential well of width $L$, such that the potential $V(x)=0$ for $0<x<L$ and $V(x)=\infty$ elsewhere. Write the time-independent Schrödinger equation for the particle inside the well and solve it to obtain the allowed energy eigenvalues. Write the normalized eigenfunctions corresponding to these energy levels.}

\vspace{0.3cm}

Inside the well, the potential is

\[
V(x)=0, \qquad 0<x<L
\]

The time-independent Schrödinger equation is

\[
-\frac{\hbar^2}{2m}\frac{d^2\psi}{dx^2}+V(x)\psi=E\psi
\]

Since $V(x)=0$ inside the well,

\[
-\frac{\hbar^2}{2m}\frac{d^2\psi}{dx^2}=E\psi
\]

Rearranging,

\[
\frac{d^2\psi}{dx^2} + \frac{2mE}{\hbar^2}\psi = 0
\]

Let

\[
k^2 = \frac{2mE}{\hbar^2}
\]

Then the equation becomes

\[
\frac{d^2\psi}{dx^2}+k^2\psi=0
\]

The general solution is

\[
\psi(x)=A\sin kx + B\cos kx
\]

The boundary conditions for an infinite potential well are

\[
\psi(0)=0, \qquad \psi(L)=0
\]

Applying $\psi(0)=0$,

\[
\psi(0)=B=0
\]

Thus,

\[
\psi(x)=A\sin kx
\]

Applying $\psi(L)=0$,

\[
A\sin kL=0
\]

For a non-trivial solution ($A\neq 0$),

\[
\sin kL=0
\]

\[
kL=n\pi, \qquad n=1,2,3,\dots
\]

\[
k=\frac{n\pi}{L}
\]

Using $k^2=\dfrac{2mE}{\hbar^2}$,

\[
\frac{2mE_n}{\hbar^2}=\frac{n^2\pi^2}{L^2}
\]

\[
E_n=\frac{n^2\pi^2\hbar^2}{2mL^2}
\]

Normalization condition:

\[
\int_0^L |\psi(x)|^2 dx =1
\]

\[
\int_0^L A^2\sin^2\left(\frac{n\pi x}{L}\right)dx=1
\]

Using

\[
\int_0^L \sin^2\left(\frac{n\pi x}{L}\right)dx=\frac{L}{2}
\]

\[
A^2\frac{L}{2}=1
\]

\[
A=\sqrt{\frac{2}{L}}
\]

Hence the normalized eigenfunctions are

\[
\psi_n(x)=\sqrt{\frac{2}{L}}\sin\left(\frac{n\pi x}{L}\right)
\]

\[
\textcolor{brown}{\mathbf{E_n=\dfrac{n^2\pi^2\hbar^2}{2mL^2}, \qquad 
\psi_n(x)=\sqrt{\dfrac{2}{L}}\sin\left(\dfrac{n\pi x}{L}\right), \quad n=1,2,3,\dots}}
\]


\textbf{5. Explain the Heisenberg gamma-ray microscope thought experiment in detail and discuss how it leads to the Heisenberg uncertainty principle, highlighting the limitations imposed on the simultaneous measurement of an electron's position and momentum.}

\vspace{0.3cm}

Consider an electron whose position is to be measured using a microscope. 
To observe the electron, it must be illuminated by radiation. 
Suppose high-frequency $\gamma$-rays of wavelength $\lambda$ are used in order to obtain high resolving power.

According to the resolving power of a microscope, the uncertainty in position is approximately

\[
\Delta x \approx \frac{\lambda}{2\sin\theta}
\]

where $\theta$ is the semi-angle of collection of the microscope objective.

To reduce $\Delta x$, one must use radiation of very small wavelength $\lambda$. 
For $\gamma$-rays,

\[
\lambda = \frac{h}{p_\gamma}
\]

where $p_\gamma$ is the momentum of the incident photon.

When the $\gamma$-ray photon strikes the electron, Compton scattering occurs. 
During this interaction, the photon transfers momentum to the electron. 
The momentum of the photon is

\[
p_\gamma = \frac{h}{\lambda}
\]

Because the scattered photon can emerge within an angle $\pm \theta$, the uncertainty in the $x$-component of the photon momentum is of order

\[
\Delta p_x \approx p_\gamma \sin\theta
\]

Since momentum is conserved, this uncertainty is transferred to the electron. Therefore,

\[
\Delta p_x \approx \frac{h}{\lambda} \sin\theta
\]

Multiplying the uncertainties in position and momentum,

\[
\Delta x \, \Delta p_x
\approx
\left( \frac{\lambda}{2\sin\theta} \right)
\left( \frac{h}{\lambda} \sin\theta \right)
\]

\[
\Delta x \, \Delta p_x
\approx
\frac{h}{2}
\]

Since $h = 2\pi\hbar$,

\[
\Delta x \, \Delta p_x \approx \pi \hbar
\]

A more rigorous quantum mechanical treatment gives

\[
\Delta x \, \Delta p_x \geq \frac{\hbar}{2}
\]

This is the Heisenberg uncertainty principle.

\vspace{0.4cm}

The thought experiment shows that in order to measure the position accurately (small $\lambda$), the photon must have large momentum. 
This large momentum disturbs the electron significantly, increasing uncertainty in its momentum. 

Conversely, using radiation of large wavelength reduces momentum disturbance but increases uncertainty in position.

Thus, the act of measurement itself introduces a fundamental limitation. 
It is not due to experimental imperfections but is intrinsic to the wave–particle duality of matter.

\vspace{0.4cm}

\[
\textcolor{brown}{\mathbf{\Delta x \, \Delta p_x \geq \dfrac{\hbar}{2}}}
\]

\[
\textcolor{brown}{\mathbf{\text{Hence, the simultaneous precise measurement of the position and momentum of an electron is fundamentally limited by the Heisenberg uncertainty principle.}}}
\]

\textbf{6. Describe Fermi theory of Beta Decay and explain why neutrinos are needed in the hypothesis.}

\vspace{0.3cm}

\textbf{Fermi Theory of Beta Decay}

Beta decay is a nuclear process in which a nucleus emits an electron (or positron) and transforms into a different nucleus. 
For $\beta^{-}$ decay, the reaction is

\[
n \rightarrow p + e^{-}
\]

However, experimentally the emitted electrons were observed to have a continuous energy spectrum rather than a discrete energy value. 
This created a serious problem because if only two particles were involved in the decay, conservation of energy and momentum would require a fixed (monoenergetic) electron energy.

To explain this, \textbf{Enrico Fermi} proposed in 1934 that beta decay is a weak interaction process involving four fermions. 
The complete $\beta^{-}$ decay process is

\[
n \rightarrow p + e^{-} + \bar{\nu}
\]

where $\bar{\nu}$ is the antineutrino.

Similarly, for $\beta^{+}$ decay,

\[
p \rightarrow n + e^{+} + \nu
\]

Fermi assumed that beta decay occurs due to a point-like interaction between four fermions (neutron, proton, electron, neutrino). 
The transition probability per unit time is given by

\[
\lambda = \frac{2\pi}{\hbar} |M_{fi}|^2 \rho(E)
\]

where $M_{fi}$ is the transition matrix element and $\rho(E)$ is the density of final states.

The interaction Hamiltonian density in Fermi's original theory is

\[
H_{int} = G_F \, \bar{\psi}_p \gamma^\mu \psi_n \, 
\bar{\psi}_e \gamma_\mu \psi_\nu
\]

where $G_F$ is the Fermi coupling constant. 
This describes a weak interaction of short range.

\vspace{0.4cm}

\textbf{Why Neutrinos are Needed}

Before neutrino hypothesis, the process was assumed as

\[
n \rightarrow p + e^{-}
\]

In such a two-body decay, conservation of energy and momentum would require the electron to have a single fixed energy.

However, experimentally:

\[
0 \leq E_e \leq E_{max}
\]

The electron energy spectrum is continuous.

This contradiction implied that another particle must be sharing the energy and momentum. 
Wolfgang Pauli proposed a neutral, very light particle — the neutrino — with:

\[
\text{Charge} = 0, \qquad \text{Spin} = \frac{1}{2}, \qquad \text{Very small mass}
\]

Including the neutrino,

\[
n \rightarrow p + e^{-} + \bar{\nu}
\]

Now energy conservation becomes

\[
E_n = E_p + E_e + E_{\bar{\nu}}
\]

Momentum conservation becomes

\[
\vec{p}_n = \vec{p}_p + \vec{p}_e + \vec{p}_{\bar{\nu}}
\]

Because the energy is shared between three particles, the electron can have a continuous range of energies, consistent with experimental observation.

The neutrino also preserves conservation of:

\[
\text{Angular momentum}
\]

since neutron and proton both have spin $\frac{1}{2}$, and the electron also has spin $\frac{1}{2}$. 
Without the neutrino, spin conservation would be violated in certain transitions.

\vspace{0.4cm}

\textcolor{brown}{\textbf{Fermi theory describes beta decay as a weak four-fermion interaction, and the neutrino is required to ensure conservation of energy, momentum, and angular momentum, thereby explaining the continuous beta spectrum.}}

\textbf{7. Explain Gamow's theory of alpha decay. Describe the quantum mechanical tunneling mechanism involved in alpha emission from a nucleus and illustrate the process with the help of a potential energy barrier diagram (graph), explaining how the alpha particle penetrates the Coulomb barrier despite having insufficient classical energy.}

\vspace{0.3cm}

In $\alpha$-decay, a heavy nucleus emits an $\alpha$-particle ($^4_2\text{He}$ nucleus) and transforms into a daughter nucleus:

\[
^{A}_{Z}X \rightarrow ^{A-4}_{Z-2}Y + \alpha + Q
\]

where $Q$ is the decay energy.

\vspace{0.4cm}

\textbf{Basic Idea of Gamow's Theory}

According to classical physics, the $\alpha$-particle confined inside the nucleus cannot escape because it does not possess enough kinetic energy to overcome the Coulomb repulsion barrier outside the nucleus.

Inside the nucleus ($r<R$), the $\alpha$-particle experiences a strong attractive nuclear potential, which can be approximated as a deep potential well.

For $r>R$, the nuclear force vanishes and the $\alpha$-particle experiences Coulomb repulsion:

\[
V(r) = \frac{1}{4\pi\varepsilon_0}\frac{2Ze^2}{r}
\]

Thus, the total potential energy has the following form:

- Deep negative well inside the nucleus.
- High positive Coulomb barrier outside.
- Potential decreases gradually as $1/r$ for large $r$.

\vspace{0.4cm}

\textbf{Quantum Mechanical Tunneling}

Although the $\alpha$-particle has total energy $E$ such that

\[
E < V_{\text{max}}
\]

quantum mechanics allows a finite probability of penetrating the barrier.

In the classically forbidden region ($R < r < b$), where $V(r) > E$, the Schrödinger equation gives exponentially decaying solutions:

\[
\psi(r) \propto e^{-\kappa r}
\]

where

\[
\kappa = \sqrt{\frac{2m(V(r)-E)}{\hbar^2}}
\]

The transmission probability (tunneling probability) is approximately

\[
T \approx e^{-2 \int_R^b \kappa(r)\, dr}
\]

This exponential dependence explains:

- Extremely long half-lives for small changes in energy.
- The Geiger–Nuttall law relating decay constant to $\alpha$ energy.

\vspace{0.4cm}

\textbf{Potential Energy Barrier Diagram (Schematic)}

\begin{center}
\begin{tikzpicture}[scale=1.0]
\draw[->] (0,0) -- (8,0) node[right] {$r$};
\draw[->] (0,0) -- (0,4) node[above] {$V(r)$};

\draw[thick] (0.5,1) -- (2,1);
\draw[thick] (2,1) .. controls (3,3.5) and (4,3.5) .. (5,2.5);
\draw[thick] (5,2.5) .. controls (6,1.8) and (7,1.2) .. (8,1);

\draw[dashed] (0,1.5) -- (8,1.5) node[right] {$E$};

\node at (1.2,0.6) {$R$};
\node at (5,0.6) {$b$};
\end{tikzpicture}
\end{center}

The horizontal dashed line represents the total energy $E$ of the $\alpha$-particle. 
Between $R$ and $b$, the region is classically forbidden. 
However, due to quantum tunneling, the wave function does not vanish but decreases exponentially, allowing a finite probability of escape.

\vspace{0.4cm}

\textbf{Conclusion of Gamow's Theory}

Gamow explained $\alpha$-decay as a quantum tunneling phenomenon where the $\alpha$-particle pre-exists inside the nucleus and escapes by penetrating the Coulomb barrier, even though its energy is less than the barrier height.

The decay constant is proportional to the tunneling probability:

\[
\lambda \propto e^{-2 \int_R^b \kappa(r)\, dr}
\]

which accounts for the observed half-lives.

\vspace{0.5cm}

\[
\textcolor{brown}{\mathbf{\alpha\text{-decay occurs due to quantum mechanical tunneling through the Coulomb barrier,}}}
\]

\[
\textcolor{brown}{\mathbf{\text{even when } E < V_{\text{max}}, \text{ with } \lambda \propto e^{-2 \int_R^b \kappa(r)\, dr}.}}
\]



\textbf{8. Derive the time-dependent Schrödinger equation starting from the concept of matter waves and hence obtain the time-independent Schrödinger equation by the method of separation of variables.}

\vspace{0.3cm}

According to de Broglie, a particle of momentum $p$ is associated with a matter wave of wavelength

\[
\lambda = \frac{h}{p}
\]

The corresponding wave number is

\[
k = \frac{2\pi}{\lambda} = \frac{p}{\hbar}
\]

The energy of the particle is related to angular frequency $\omega$ as

\[
E = \hbar \omega
\]

A plane matter wave may therefore be written as

\[
\Psi(x,t) = A e^{i(kx-\omega t)}
\]

Differentiating with respect to time,

\[
\frac{\partial \Psi}{\partial t} = -i\omega \Psi
\]

\[
i\hbar \frac{\partial \Psi}{\partial t} = \hbar \omega \Psi = E\Psi
\]

Differentiating twice with respect to $x$,

\[
\frac{\partial^2 \Psi}{\partial x^2} = -k^2 \Psi
\]

\[
-\hbar^2 \frac{\partial^2 \Psi}{\partial x^2} = \hbar^2 k^2 \Psi = p^2 \Psi
\]

Using the classical relation for total energy,

\[
E = \frac{p^2}{2m} + V(x,t)
\]

Substituting the operator forms,

\[
i\hbar \frac{\partial \Psi}{\partial t}
=
\frac{1}{2m}
\left(
-\hbar^2 \frac{\partial^2}{\partial x^2}
\right)\Psi
+
V(x,t)\Psi
\]

Thus, the time-dependent Schrödinger equation in one dimension is

\[
i\hbar \frac{\partial \Psi(x,t)}{\partial t}
=
-\frac{\hbar^2}{2m}
\frac{\partial^2 \Psi(x,t)}{\partial x^2}
+
V(x,t)\Psi(x,t)
\]

In three dimensions,

\[
i\hbar \frac{\partial \Psi}{\partial t}
=
-\frac{\hbar^2}{2m}\nabla^2 \Psi
+
V\Psi
\]

\vspace{0.4cm}

\textbf{Derivation of the Time-Independent Schrödinger Equation}

If the potential is independent of time, $V = V(x)$, we assume a separable solution

\[
\Psi(x,t) = \psi(x)T(t)
\]

Substituting into the time-dependent equation,

\[
i\hbar \psi(x)\frac{dT}{dt}
=
-\frac{\hbar^2}{2m}T(t)\frac{d^2\psi}{dx^2}
+
V(x)\psi(x)T(t)
\]

Dividing by $\psi(x)T(t)$,

\[
i\hbar \frac{1}{T}\frac{dT}{dt}
=
-\frac{\hbar^2}{2m}\frac{1}{\psi}\frac{d^2\psi}{dx^2}
+
V(x)
\]

The left-hand side depends only on $t$, and the right-hand side only on $x$, therefore both must be equal to a constant $E$.

From the time part,

\[
i\hbar \frac{1}{T}\frac{dT}{dt} = E
\]

\[
T(t) = e^{-iEt/\hbar}
\]

From the spatial part,

\[
-\frac{\hbar^2}{2m}\frac{d^2\psi}{dx^2}
+
V(x)\psi
=
E\psi
\]

This is the time-independent Schrödinger equation.

\vspace{0.5cm}

\[
\textcolor{brown}{\mathbf{i\hbar \frac{\partial \Psi}{\partial t}
=
-\frac{\hbar^2}{2m}\nabla^2 \Psi + V\Psi}}
\]

\[
\textcolor{brown}{\mathbf{-\frac{\hbar^2}{2m}\nabla^2 \psi + V\psi = E\psi}}
\]


\textbf{9. Write short notes on any two of the following:}

\vspace{0.3cm}

%%%%%%%%%%%%%%%%%%%%%%%%%%%%%%%%%%%%%%%%%%%%%%%%%%%%%%%%%%%%
\textbf{(a) Define half-life and derive the expression for the half-life of a radioactive substance.}

The rate of radioactive decay is proportional to the number of undecayed nuclei present at any time:

\[
\frac{dN}{dt} = -\lambda N
\]

where $\lambda$ is the decay constant.

Separating variables,

\[
\frac{dN}{N} = -\lambda dt
\]

Integrating from $N_0$ at $t=0$ to $N$ at time $t$,

\[
\int_{N_0}^{N} \frac{dN}{N}
=
-\lambda \int_{0}^{t} dt
\]

\[
\ln \left( \frac{N}{N_0} \right) = -\lambda t
\]

\[
N = N_0 e^{-\lambda t}
\]

Half-life $T_{1/2}$ is defined as the time required for the number of nuclei to reduce to half of its initial value:

\[
N = \frac{N_0}{2}
\]

Substituting,

\[
\frac{N_0}{2} = N_0 e^{-\lambda T_{1/2}}
\]

\[
\frac{1}{2} = e^{-\lambda T_{1/2}}
\]

Taking logarithm,

\[
\ln 2 = \lambda T_{1/2}
\]

\[
T_{1/2} = \frac{\ln 2}{\lambda}
\]

\[
\textcolor{brown}{\mathbf{T_{1/2} = \dfrac{\ln 2}{\lambda} = \dfrac{0.693}{\lambda}}}
\]

\vspace{0.8cm}

%%%%%%%%%%%%%%%%%%%%%%%%%%%%%%%%%%%%%%%%%%%%%%%%%%%%%%%%%%%%
\textbf{(b) Derive and explain the rate equations for population densities in a two-level laser system and discuss why population inversion cannot be achieved under steady-state conditions.}

Consider a two-level system with lower level $E_1$ and upper level $E_2$ having populations $N_1$ and $N_2$ respectively.

Absorption rate:

\[
\text{Rate}_{\text{abs}} = B_{12} \rho(\nu) N_1
\]

Stimulated emission rate:

\[
\text{Rate}_{\text{stim}} = B_{21} \rho(\nu) N_2
\]

Spontaneous emission rate:

\[
\text{Rate}_{\text{sp}} = A_{21} N_2
\]

The rate equation for the upper level is

\[
\frac{dN_2}{dt}
=
B_{12}\rho(\nu)N_1
-
B_{21}\rho(\nu)N_2
-
A_{21}N_2
\]

Using $N_1 + N_2 = N$,

Under steady-state condition,

\[
\frac{dN_2}{dt}=0
\]

\[
B_{12}\rho(\nu)N_1
=
B_{21}\rho(\nu)N_2
+
A_{21}N_2
\]

\[
\frac{N_2}{N_1}
=
\frac{B_{12}\rho(\nu)}
{B_{21}\rho(\nu) + A_{21}}
\]

Even for very large radiation density $\rho(\nu)$,

\[
\frac{N_2}{N_1}
\rightarrow
\frac{B_{12}}{B_{21}}
\]

From Einstein relations,

\[
B_{12} = B_{21}
\]

Therefore,

\[
\frac{N_2}{N_1} \leq 1
\]

Thus, population inversion ($N_2 > N_1$) cannot be achieved in a two-level system under steady-state conditions.

\[
\textcolor{brown}{\mathbf{\text{Population inversion cannot be achieved in a two-level system since } N_2 \leq N_1 \text{ under steady state.}}}
\]


%%%%%%%%%%%%%%%%%%%%%%%%%%%%%%%%%%%%%%%%%%%%%%%%%%%%%%%%%%%%
\textbf{(c) Physical interpretation of wave function}

In quantum mechanics, the state of a particle is described by a wave function 
\[
\Psi(x,t)
\]

The wave function itself has no direct physical meaning. However, according to Born’s probabilistic interpretation, the physically measurable quantity is the probability density:

\[
|\Psi(x,t)|^2 = \Psi^*(x,t)\Psi(x,t)
\]

which represents the probability per unit length of finding the particle at position $x$ at time $t$.

The probability of finding the particle between $x$ and $x+dx$ is

\[
|\Psi(x,t)|^2 dx
\]

For a physically acceptable wave function, the following conditions must be satisfied:

It must be single-valued.

It must be finite everywhere.

It must be continuous, and its first derivative must be continuous (except where the potential is infinite).

It must be normalizable:

\[
\int_{-\infty}^{\infty} |\Psi(x,t)|^2 dx = 1
\]

Expectation value of any observable $\hat{O}$ is given by

\[
\langle O \rangle = \int \Psi^* \hat{O} \Psi \, dx
\]

Thus, the wave function contains complete information about the system through probability amplitudes.

\[
\textcolor{brown}{\textbf{The wave function itself has no direct physical meaning; its modulus squared }}
\]
\[ |\Psi|^2 \text{ represents the probability density of finding the particle.}\]

\vspace{0.8cm}

%%%%%%%%%%%%%%%%%%%%%%%%%%%%%%%%%%%%%%%%%%%%%%%%%%%%%%%%%%%%
\textbf{(d) de-Broglie's hypothesis on wave-particle duality and explain its physical significance}

According to de Broglie, not only radiation but also material particles exhibit wave-like properties.

For a particle of momentum $p$,

\[
\lambda = \frac{h}{p}
\]

This is known as the de Broglie wavelength.

For a particle of mass $m$ moving with velocity $v$,

\[
\lambda = \frac{h}{mv}
\]

Using the energy relation,

\[
E = h\nu
\]

and

\[
p = \hbar k
\]

de Broglie unified the concepts of particle and wave nature.

Physical significance:

It explains electron diffraction experiments such as Davisson–Germer experiment.

It leads to quantization condition in Bohr’s atomic model:

\[
2\pi r = n\lambda
\]

It forms the basis of wave mechanics and Schrödinger equation.

It establishes the fundamental principle of wave-particle duality.

\[
\textcolor{brown}{\mathbf{\lambda = \dfrac{h}{p}}}
\]

\[
\textcolor{brown}{\textbf{Matter particles exhibit wave nature, establishing the principle of wave-particle duality.}}
\]










%------------Body-Ends--------------------
%\bibliography{ref.bib}
%\bibliographystyle{IEEEtran}
\end{document}
